\documentclass[11pt]{article}
\usepackage[margin=1in]{geometry}
\usepackage{amsmath,amssymb,hyperref}
\newcommand{\T}{^{\mathsf T}}
\newcommand{\sym}{\operatorname{sym}}
\newcommand{\ip}[2]{\langle #1,#2\rangle}
\title{Working equations for the RHF DLPNO-CCSD(T) nuclear gradient}
\author{Quantum Chemistry Grand Battleplan}
\date{Occupied-space projection and frozen-core response}
\begin{document}
\maketitle
This is a cumulative implementation record. Equation labels identify new
derivative rules. The forward energy, amplitudes, local transformations and
published Lambda equations retain the definitions of Jiang et al.\ (2024)
and T\'oth et al.\ (2026). This edition covers propagation from the raw PAO
and occupied-operator sources to the AO and orbital inputs. All differentials
hold pair lists, ranks, and the identities of the frozen orbitals fixed.
\section{Indices and derivative convention}
As in Jiang, $\mu,\nu,\lambda,\sigma$ label AOs, $\tilde\nu$ a PAO
identified by its parent AO $\nu$, $i,j,k,l$ active occupied LMOs, and $a,b$
virtual orbitals. We reserve $i_{\rm c}$ for a frozen occupied orbital.
A bar denotes a coefficient of the stationary correlation Lagrangian
differential. For real matrices we use
\begin{equation}\label{eq:pairing}
 \ip{X}{Y}=\sum_{pq}X_{pq}Y_{pq},\qquad
 \sym X=(X+X\T)/2,\qquad d\mathcal L=\sum_X\ip{\bar X}{dX}.
\end{equation}
In particular, off-diagonal elements of a symmetric input occur twice in
this pairing. Let $L_{\mu i}$ be active localized coefficients, $C^{\rm c}_{\mu i_{\rm c}}$
the reference frozen coefficients, and $O=[L\;C^{\rm c}]$. $S,F$ denote
the AO overlap and RHF Fock matrices.
\section{Unnormalized projected AO}
For a spatial RHF occupied projector each occupied column occurs once.
Define $P=OO\T$. In the AO coefficient representation the raw PAO and its
variation are
\begin{align}
 B_{\mu\tilde\nu}
 &=\delta_{\mu\nu}-\sum_{p\in\mathrm{occ}}\sum_\lambda
 O_{\mu p}O_{\lambda p}S_{\lambda\nu},\label{eq:raw}\\
 dB&=-(dO\,O\T+O\,dO\T)S-P\,dS.\label{eq:db}
\end{align}
For $O\T S O=I$, $O\T S B=0$. The overlap in the rightmost position
matters because $B$ is generally nonsymmetric. If $b$ is one raw column
and $x=b/n$, $n^2=b\T S b>0$, a normalized-column source $g$ gives
\begin{equation}\label{eq:normal}
 \bar b=g/n-Sb(g\T b)/n^3,\qquad
 \bar S_{\rm norm}=-bb\T(g\T b)/(2n^3).
\end{equation}
Apply the actual forward zero-column convention before this rule.
When normalization is differentiated using the diagonal of $B\T S B$,
its source enters the metric block below and must be counted once.
\section{Inputs from the pair and triplet contractions}
Let $U,W,V,H,D,R,Z$ denote the accumulated sources for, respectively,
$L\T F L$, $L\T F B$, $B\T F B$, $B\T S B$, $B$, $S$, and $SL$.
The last two include the local occupied density-fitting refits.
These definitions specify their normalization by the following scalar
differential, with the sources held fixed:
\begin{align}
 d\mathcal L_{\rm loc}
 &=\ip U{d(L\T F L)}+\ip W{d(L\T F B)}
   +\ip V{d(B\T F B)}+\ip H{d(B\T S B)}\notag\\
 &\quad+\ip D{dB}+\ip R{dS}+\ip Z{d(SL)}.\label{eq:sources}
\end{align}
Thus the total source on the raw PAO coefficients at fixed $L,F,S$ is
\begin{equation}\label{eq:g}
 G=D+FLW+FB(V+V\T)+SB(H+H\T).
\end{equation}
\section{Occupied, frozen, Fock and metric derivatives}
Substitution of Eq.~\eqref{eq:db} in Eq.~\eqref{eq:sources} gives
\begin{align}
 \bar L&=-GSL-SG\T L+FL(U+U\T)+FBW\T+SZ,\label{eq:l}\\
 \bar C^{\rm c}&=-GSC^{\rm c}-SG\T C^{\rm c},\label{eq:c}\\
 \bar F&=\sym(LUL\T+LWB\T+BVB\T),\label{eq:f}\\
 \bar S&=\sym(BHB\T+R+ZL\T-PG).\label{eq:s}
\end{align}
For example the individual frozen-column equation reads
\begin{equation}\label{eq:coreindex}
 \bar C^{\rm c}_{\mu i_{\rm c}}
 =-\sum_{\nu\lambda}
 \left(G_{\mu\tilde\nu}S_{\nu\lambda}
 +S_{\mu\nu}G_{\lambda\tilde\nu}\right)
 C^{\rm c}_{\lambda i_{\rm c}}.
\end{equation}
With no frozen orbitals this is an AO-by-zero block. For an infinitesimal
rotation $dC^{\rm c}_{\mu i_{\rm c}}=C_{\mu a}\,d\kappa_{a i_{\rm c}}$
at fixed active LMOs and AO inputs, the source is
\begin{equation}\label{eq:kappa}
 \frac{\partial\mathcal L_{\rm loc}}{\partial\kappa_{a i_{\rm c}}}
 =\sum_\mu C_{\mu a}\bar C^{\rm c}_{\mu i_{\rm c}}.
\end{equation}
This orbital source is an input to the constrained RHF response problem.
A rotation within the frozen block, $dC^{\rm c}=C^{\rm c}K$ with $K=-K\T$,
leaves its projector invariant by Eq.~\eqref{eq:db}.
\section{Implementation and verification}
\texttt{occupied\_projection.cc} applies Eqs.~\eqref{eq:l}--\eqref{eq:s}
to a supplied column panel. For example $BX=X-O(O\T SX)$ and
$B\T X=X-SO(O\T X)$. The sources $U,W,V,H,D,R,Z$ correspond in order
to \texttt{ccsd\_foo\_response\_}, the raw \texttt{fov}, \texttt{fvv},
\texttt{overlap}, \texttt{coefficients}, \texttt{df\_refit\_overlap\_bar\_},
and \texttt{df\_sclmo\_bar\_}. The debug probe perturbs one independent
input at fixed geometry and reevaluates all seven products of
Eq.~\eqref{eq:sources}. The separate canonical test rotates the frozen
orbital toward a virtual direction and evaluates the stationary DF-CCSD(T)
correlation Lagrangian. The final nuclear derivative additionally requires
the localization and RHF constraints and the remaining explicit integral
derivatives.
\end{document}
